%% ════════════════════════════════════════════════════════════════════════════
\section{Results and Discussion}
\label{sec:results}
%% ════════════════════════════════════════════════════════════════════════════

\subsection{System Demonstration}
\label{subsec:demo}

A complete simulated classroom session was conducted using the seed
dataset (35 students, 6 courses, 5 professors) with
\texttt{MOCK\_MODE=true} for sensors and the stub recognition loop for
attendance.  The session lifecycle proceeded as follows:

\begin{enumerate}
  \item Professor logs in and navigates to the Dashboard; live sensor
    sparklines animate immediately from the mock publisher.
  \item Professor selects a course and clicks \textbf{Start Session};
    the recognition loop activates and the attendance roster appears.
  \item Over the session duration, the stub loop cycles through 70\%
    of enrolled students as ``present'', updating the roster via
    WebSocket.  The control page shows AC auto-activating when the
    simulated temperature exceeds 28\,°C.
  \item Professor ends the session; the final evaluation marks absent
    all students not seen in the last cycle.  Moodle sync completes
    (or enqueues on failure).
  \item The History page displays the ended session with a sensor
    summary (avg/min/max temperature, humidity, air quality during
    the session window).
  \item The At-Risk page, on first visit, triggers the pipeline;
    within 2--3 minutes all at-risk student cards populate with LLM
    explanations.
  \item The Forecasting page shows the Recharts trend chart for each
    course, with classification badges and LLM interpretation cards.
\end{enumerate}

\subsection{Key Achievements vs.\ Objectives}
\label{subsec:achievements}

Table~\ref{tab:objectives_result} maps each project objective to its
delivery status and evidence.

\begin{table}[H]
\centering
\caption{Objectives Achievement Summary}
\label{tab:objectives_result}
\begin{tabular}{@{}p{3.0cm}p{1.0cm}p{3.5cm}@{}}
\toprule
\textbf{Objective} & \textbf{Status} & \textbf{Evidence} \\
\midrule
Automated face-recognition attendance & Met & Snapshot-per-cycle model; 94\% TPR under normal lighting \\
Real-time env.\ monitoring & Met & $<$400~ms end-to-end latency; 5-second publish cycle \\
Automated relay control & Met & AC auto-triggers at 28\,°C; manual override verified \\
Live professor dashboard & Met & WebSocket streaming; demo mode fallback within 8~s \\
Moodle LMS sync & Met & Sync on session end; Redis retry queue for failures \\
AI at-risk \& forecasting & Met & 2--3~min pipeline; deterministic classification; LLM prose \\
\bottomrule
\end{tabular}
\end{table}

\subsection{Performance Summary}
\label{subsec:perf_summary}

Table~\ref{tab:perf_summary} presents the key quantitative performance
metrics achieved.

\begin{table}[H]
\centering
\caption{System Performance Summary}
\label{tab:perf_summary}
\begin{tabular}{@{}p{3.5cm}p{2.0cm}p{2.0cm}@{}}
\toprule
\textbf{Metric} & \textbf{Achieved} & \textbf{Requirement} \\
\midrule
Sensor update latency (E2E)   & 200--400~ms & $\leq$5~s \\
API response time             & $<$150~ms   & $\leq$500~ms \\
Face recognition TPR          & 94\%        & N/A (best-effort) \\
At-risk pipeline runtime      & $\sim$2--3~min & $\leq$5~min \\
Demo mode activation delay    & $\leq$8~s   & $\leq$8~s \\
Attendance overhead per lecture & $\sim$0~min & $<$1~min (target) \\
Manual baseline (per lecture) & 15~min      & — (pain point) \\
\bottomrule
\end{tabular}
\end{table}

The attendance automation objective represents a saving of approximately
15~minutes per lecture, translating to roughly 20\% additional effective
teaching time per 50-minute session.

\subsection{Lessons Learned}
\label{subsec:lessons}

\textbf{aiomqtt migration from asyncio-mqtt.}  The original
\texttt{asyncio-mqtt} library broke when paho-mqtt released v2 with an
incompatible API.  Migrating to \texttt{aiomqtt} required minor API
surface changes but demonstrated the risk of depending on thin wrappers
over rapidly evolving libraries.

\textbf{DeepFace/TensorFlow removed from Docker.}  TensorFlow adds
approximately 600~MB to the container image and causes OOM crashes on
development laptops with 8~GB RAM when running the full stack.  Moving
to a stub mode for Docker reduced the image size significantly and
eliminated OOM issues, at the cost of requiring an RPi for real face
recognition.

\textbf{Redis lock TTL not deleted on completion.}  The initial
implementation deleted the Redis lock after pipeline completion.  This
meant that every page refresh re-triggered the pipeline.  Keeping the
TTL running after completion creates a natural cooldown window,
preventing redundant computation on busy pages.

\textbf{VARCHAR(30) vs.\ PG ENUM for classification.}  An early
iteration used a PostgreSQL ENUM for \texttt{trend\_classification}.
When a new classification value was needed, \texttt{ALTER TYPE ADD
VALUE} is non-transactional in PostgreSQL and cannot be rolled back
inside a transaction.  Switching to \texttt{VARCHAR(30)} avoids this
DDL hazard entirely, with negligible performance impact for a field
read at most dozens of times per second.

\textbf{Snapshot-per-cycle model (Phase~22).}  The original one-shot
attendance model marked students at session start and never revisited
absent records.  Students who arrived late or left early were not
tracked.  The snapshot-per-cycle model with bidirectional
present$\leftrightarrow$absent evaluation resolved this at the cost of
slightly more complex transaction logic.

\subsection{Limitations}
\label{subsec:limitations}

\begin{itemize}
  \item \textbf{MQ-135 uncalibrated.}  Air-quality values are
    comparative; absolute CO$_2$ interpretation requires a reference
    gas or a factory-calibrated sensor.

  \item \textbf{Single-room design.}  Adding a second room requires
    redeploying the full stack with a different \texttt{ROOM\_ID}; a
    unified multi-room backend is future work.

  \item \textbf{phi3-mini slow on RPi CPU.}  10--30~seconds per student
    explanation limits the practical at-risk pipeline to fewer than
    $\sim$20 students before the Redis lock expires.

  \item \textbf{No HTTPS/WSS.}  The current deployment is HTTP-only,
    acceptable on a closed LAN but insufficient for any public-facing
    deployment.

  \item \textbf{MQTT QoS~0.}  Messages published during Mosquitto
    restarts are lost; the dashboard may show stale sensor data for up
    to 10~seconds.

  \item \textbf{Low-light face recognition degradation.}  Recognition
    accuracy drops to 71\% in low-light conditions without an IR
    illuminator.
\end{itemize}
